% Settlement — the compute ledger of the Hanzo sovereign L1 on Lux.
\documentclass[11pt,a4paper]{article}
\usepackage{../shared/hanzocover}
\usepackage[utf8]{inputenc}
\usepackage[T1]{fontenc}
\usepackage{amsmath,amssymb,amsfonts,amsthm}
\usepackage{mathtools}
\usepackage{booktabs}
\usepackage{longtable}
\usepackage{tabularx}
\usepackage{hyperref}
\usepackage{algorithm}
\usepackage{algpseudocode}
\usepackage{enumitem}
\usepackage{xcolor}
\usepackage{listings}
\input{../shared/lstlang}

\hypersetup{colorlinks=true,linkcolor=black,citecolor=blue,urlcolor=blue}

\newtheorem{theorem}{Theorem}[section]
\newtheorem{lemma}[theorem]{Lemma}
\newtheorem{proposition}[theorem]{Proposition}
\newtheorem{corollary}[theorem]{Corollary}
\theoremstyle{definition}
\newtheorem{definition}{Definition}[section]
\newtheorem{invariant}{Invariant}[section]
\newtheorem{observation}{Observation}[section]
\newtheorem{parameter}{Parameter}[section]
\theoremstyle{remark}
\newtheorem{remark}{Remark}[section]

\lstset{
  basicstyle=\ttfamily\footnotesize,
  breaklines=true,
  frame=single,
  numbers=none,
  xleftmargin=1.2em,
  columns=fullflexible,
}

\newcommand{\code}[1]{\texttt{#1}}
% File paths and identifiers do not hyphenate, so a paragraph of them needs
% extra interword stretch rather than a hyphenation rule that would break
% names in misleading places.
\sloppy
\emergencystretch=3em
\newcommand{\Ord}{\omega}
\newcommand{\Claim}{c}
\newcommand{\Ops}{\mathcal{O}}
\newcommand{\Buyers}{\mathcal{B}}
\newcommand{\Admit}{\textsc{Admit}}

\title{Settlement\\[6pt]
\large The Compute Ledger of the Hanzo Sovereign L1 on Lux}

\author{
Hanzo AI Research\thanks{Corresponding author: research@hanzo.ai}\\
\textit{Hanzo Industries, Inc.\ (Techstars '17)}\\
\texttt{research@hanzo.ai}
}

\date{2026}

\begin{document}

\hanzocoverpage

\begin{abstract}
Hanzo settles paid AI compute. Buyers put credit in; operators who serve
attested work take rewards out. This document specifies the ledger that
connects those two flows, the predicate that decides whether a served-compute
claim is admitted, the bond an operator posts and the conditions under which it
is reduced, and the exact boundary between what the chain records and what stays
off it.

The chain is a sovereign L1 on Lux. No consensus, cryptography, or
validator-set machinery is introduced here: ordering and finality are Lux
Quasar, the post-quantum primitives are the Lux precompile block, the
attestation verifier is the one already compiled into that block, and the L1 is
born by \code{CreateChainTx} + \code{AddValidatorTx} + \code{ConvertNetworkTx}
against the Lux P-chain and then runs on its own validator set with an on-chain
manager. The contribution of this document is the accounting rules and the
admission predicate, not a chain.

We separate two money flows that the current implementation braids together. A
buyer's escrowed credit pays an operator for work the buyer ordered: a transfer,
conserved, bounded by the order. A protocol emission additionally rewards
attested work: a mint, bounded by a per-epoch pool and distributed pro-rata.
Reading the deployed code, the second flow is presently unbounded --- a single
genuinely attested device can compute a reward of up to $6.44 \times 10^{9}$~AI
from one claim, and can produce unboundedly many distinct claims, because the
reward magnitude is read from an attacker-chosen field of the work proof and the
replay key is derived from an attacker-chosen nonce. We give the specific
predicate and accounting changes that close this, each expressed as a
composition of primitives already present in \code{luxfi/precompile}: the
existing ML-DSA verifier, the existing TEE device binding, and the existing
spent set.

Throughout we mark what is observed in code from what is proposed, and cite the
file and symbol for every claim about current behaviour.
\end{abstract}

\tableofcontents
\newpage

%=====================================================================
\section{Scope and method}

\subsection{What this document specifies}

Five things, and nothing else:

\begin{enumerate}[label=(\arabic*),leftmargin=2.2em]
\item the ledger --- pay-per-use compute credit in, AI-work reward out;
\item operator bonds and the conditions under which a bond is reduced,
      including a missing or invalid attestation;
\item the distribution of reward for served, attested compute;
\item the admission of a served-compute claim: the attestation proof carried
      from the serving daemon or its TEE, the predicate that accepts it, and
      the settlement that follows;
\item the boundary between on-chain and off-chain state.
\end{enumerate}

\subsection{What this document does not do}

It does not design a chain. Ordering, finality, the validator wire, the
post-quantum signature suite, the threshold suite, the TEE certificate-chain
verifier, and the cross-chain message format are all Lux machinery, cited where
used and not restated. It does not set token supply or an emission schedule: the
published record disagrees with itself on both (\S\ref{sec:disagreement}), so we
specify the invariant the schedule must satisfy and leave the schedule to
governance. It does not specify the market that decides which operator serves
which order; that is the clearing problem, treated separately.

\subsection{Evidential conventions}

Every statement about current behaviour is labelled and carries a file
reference. \textbf{Observed} means the behaviour is present in code we read.
\textbf{Derived} means it follows from observed code by an argument we give.
\textbf{Proposed} means it does not exist and is put forward here.
\textbf{Assumed} means it is a security or economic premise we do not prove.
Repository roots: \code{luxfi/node}, \code{luxfi/precompile}, \code{luxfi/evm},
\code{luxfi/consensus} under \code{\textasciitilde/work/lux}; the off-chain
plane is \code{hanzoai/cloud}.

The key words MUST, MUST NOT, SHOULD, SHOULD NOT and MAY are to be interpreted
as in RFC 2119.

%=====================================================================
\section{The chain}
\label{sec:chain}

\subsection{Birth on the Lux P-chain}

\textbf{Observed.} The Lux P-chain distinguishes the construction of a network
from its promotion. \code{CreateNetworkTx} is the constructor
($\emptyset \to \text{Network}$); \code{ConvertNetworkTx} is the endomorphism
($\text{Network} \to \text{Network}$) that changes a network's security from
inherited to sovereign and re-anchors its parent
(\code{node/vms/platformvm/txs/convert\_network\_tx.go}). Its documentation is
explicit that this is the L2-to-L1 promotion.

The security configuration carried by that transaction is a value with two
orthogonal axes (\code{node/vms/platformvm/security/mode.go}):

\begin{lstlisting}[language=Go]
type Mode struct {
    RestakeParent bool
    Admission     Admission // NoOwnSet | Open | Gated
    Threshold     uint64    // minimum stake for Open admission
    Manager       Manager   // PChain | Contract
}
func (m Mode) Sovereign() bool { return m.Admission != NoOwnSet }
\end{lstlisting}

\code{SyntacticVerify} refuses a conversion whose resulting mode is not
sovereign (\code{ErrConvertMustEstablishOwnSet}), refuses one with no genesis
validators (\code{ErrConvertMustHaveValidators}), and refuses
\code{Manager == Contract} without a manager address
(\code{ErrContractManagerNeedsAddress}).

\textbf{Observed.} \code{AddValidatorTx} is the universal entry point for
adding a validator to a network; \code{AddChainValidatorTx} carries a
deprecation notice stating that under LP-018 validators join a network and never
a chain, and that ``a sovereign L1 IS a primary network at its own networkID''
(\code{add\_chain\_validator\_tx.go}). \code{CreateChainTx} creates a blockchain
on a network.

\textbf{Specification.} Hanzo is born by the sequence \code{CreateNetworkTx}
(once, if the network does not already exist), \code{CreateChainTx},
\code{AddValidatorTx}, \code{ConvertNetworkTx}, executed against the upstream
Lux P-chain, with

\[
\texttt{Mode} = \{\ \texttt{RestakeParent}: \text{false},\
\texttt{Admission}: \texttt{Open},\
\texttt{Threshold}: \tau_{\min},\
\texttt{Manager}: \texttt{Contract}\ \}
\]

and \code{ManagerChainID} equal to the Hanzo EVM chain, \code{ManagerAddress}
equal to the staking contract of \S\ref{sec:bond}. After the conversion the L1
runs on its own validator set; the Lux P-chain is used at runtime only as the
registry that mirrors that set.

\subsection{Network identity}

\textbf{Convention.} A sovereign L1 runs its own primary network whose network
identifier equals its EVM chain identifier, one identifier per environment. This
makes staking-key derivation paths globally unique and removes wallet ambiguity.
The node's own source states the first half of this directly: ``a sovereign L1
IS a primary network at its own networkID''.

\textbf{Observed.} Two identifiers are already assigned and in use: mainnet
$36963$ (\code{hanzo-genesis.json} \code{config.chainId}; deployment records in
\code{hanzo/node/contracts/lib/standard/DEPLOYMENTS.md}) and testnet $36962$
(HIP-0024). \textbf{Proposed.} Devnet and localnet identifiers are unassigned;
they MUST be allocated through the chain registry (HIP-1020) before use, and
this document does not assign them.

\subsection{What finality authorises}

\textbf{Observed.} Lux finality is a ladder, not a boolean, and each rung
authorises a strictly larger class of action
(\code{consensus/engine/chain/finality.go}):

\begin{center}
\begin{tabular}{@{}lll@{}}
\toprule
Rung & Condition & Authorises \\
\midrule
Photon  & vote in flight & nothing \\
Wave    & forming preference & nothing \\
Nova    & $\beta$-ignition, majority stake & local execution; reorgable \\
Quasar  & $>\!2/3$ stake certificate & export: bridges, cross-chain, settlement \\
Horizon & post-quantum seal of that certificate & irreversible settlement \\
\bottomrule
\end{tabular}
\end{center}

The predicates are \code{AuthorizesLocalExecution() \{ return f >= Nova \}},
\code{AuthorizesExport() \{ return f >= Quasar \}} and
\code{AuthorizesIrreversibleSettlement() \{ return f >= Horizon \}}. The stake
floors have single definitions, \code{config.TwoThirdsStakeFloor} and
\code{config.HalfStakeFloor}, that the certificate verifier enforces.

\textbf{Specification.} Each money movement in this document is bound to a rung
in \S\ref{sec:finality}. The rule is that a movement which another party can act
on irreversibly requires the rung that authorises irreversibility, and no
movement is bound to a rung looser than the consequence it enables.

\subsection{What is reused, and where it lives}

Packages below are in \code{luxfi/precompile} unless the row names another
repository; slot addresses are given by their significant prefix.

\begin{center}
\small
\begin{tabular}{@{}lll@{}}
\toprule
Concern & Mechanism & Location \\
\midrule
Ordering, finality & Quasar & \code{consensus}, \code{engine/chain} \\
Signature over a claim & ML-DSA (FIPS 204) & \code{ai}, slot \code{0x012202} \\
TEE quote verification & X.509 chain to root & \code{ai/tee\_attest.go} \\
Attestation ABI & selectors \code{0x01}--\code{0x05} & \code{attestation}, \code{0x0300\ldots01} \\
Replay prevention & spent set & \code{ai}, \code{IsSpent}/\code{MarkSpent} \\
Model commitment & versioned registry & \code{modelregistry}, \code{0x0300\ldots02} \\
Deterministic re-execution & integer transformer & \code{inference}, \code{0x0300\ldots03} \\
Market state & provider/job records & \code{compute}, \code{0x0300\ldots10} \\
Validator stake read & signed ballot & \code{evm}: \code{stakeweight}, \code{0x0200\ldots06} \\
Cross-chain message & Warp & \code{evm}: \code{warp}, \code{0x0200\ldots05} \\
Off-chain state commitment & anchor roots & \code{anchor}, \code{0x0700\ldots10} \\
\bottomrule
\end{tabular}
\end{center}

\textbf{Observed.} All of these except \code{stakeweight} are activated on Lux
mainnet: \code{genesis/configs/mainnet/upgrade.json} lists 49
\code{precompileUpgrades} keys at \code{blockTimestamp 1766708400}, including
\code{aiMiningConfig}, \code{attestationConfig}, \code{computeMarketConfig},
\code{anchorConfig}, \code{warpConfig}, \code{mldsaVerify} and
\code{rewardManagerConfig}. \code{stakeWeightConfig} is implemented
(\code{evm/precompile/contracts/stakeweight}) but does not appear in that file,
so it is not active on mainnet today; \S\ref{sec:bond} depends on its
activation.

%=====================================================================
\section{Objects}
\label{sec:objects}

We fix notation before stating any rule. $\Buyers$ is the set of buyer
accounts, $\Ops$ the set of operator accounts. All amounts are integers in the
smallest unit of AI ($10^{-18}$~AI); no rule in this document uses a
non-integer.

\begin{definition}[Model commitment]
$\mu \in \{0,1\}^{256}$ is the identifier of a model version in the on-chain
registry. \textbf{Observed:} \code{precompile/modelregistry} stores, per model
name, a weight commitment and a version under admin control
(\code{weightKey}, \code{verKey}, \code{IsAdmin}). $\mu$ names the pair; it
does not prove that the named weights produced any particular output.
\end{definition}

\begin{definition}[Order]
$\Ord = (\iota, b, \mu, \pi, \bar m, \delta, o)$, where $\iota$ is the order
identifier, $b \in \Buyers$ the buyer, $\mu$ the model commitment, $\pi$ the
price cap in the smallest unit, $\bar m$ the meter cap (the maximum quantity of
work the buyer agrees to pay for, in the meter unit of \S\ref{sec:meter}),
$\delta$ the deadline as a block height, and $o \in \Ops \cup \{\bot\}$ the
assigned operator ($\bot$ = open to any bonded operator).
\end{definition}

\begin{definition}[Work proof]
The byte layout the AI precompile already parses
(\code{precompile/ai/ai\_mining.go}, \code{ParseWorkProof}):
\begin{align*}
&\text{deviceID}(32)\ \|\ \text{nonce}(32)\ \|\ \text{timestamp}(\text{u64})\\
&\qquad \|\ \text{privacy}(\text{u16})\ \|\ \text{computeMinutes}(\text{u32})\ \|\ \text{teeQuote}(\ast)
\end{align*}
\end{definition}

\begin{definition}[Attestation envelope]
The envelope appended to a work proof
(\code{precompile/ai/cross\_chain\_mining.go}, \code{verifyDeviceBinding}):
\begin{align*}
\text{envelope} &= \text{sigLen}(\text{u32})\ \|\ \text{teeSig}\ \|\ \text{receipt}\\
\text{receipt} &= \underbrace{\text{deviceID}(32)\|\text{ts}(8)\|\text{nonce}(8)}_{\text{report, 48 B}}\ \|\ \text{certChainDER}
\end{align*}
\end{definition}

\begin{definition}[Claim]
$\Claim = (\iota, o, P, \sigma, \alpha, y, m, \lambda)$: the order identifier,
the operator, the ML-DSA public key $P$, the ML-DSA signature $\sigma$ over the
canonical claim bytes, the attestation envelope $\alpha$, the output commitment
$y$, the meter reading $m$, and the privacy level $\lambda$.
\end{definition}

\begin{definition}[Privacy multiplier]
$\phi: \lambda \mapsto \mathbb{N}$, in basis points. \textbf{Observed:}
$\phi$ is \code{privacyMultipliers} in \code{precompile/ai/ai\_mining.go}, with
$\phi(1) = 2500$ (public), $\phi(2) = 5000$ (private), $\phi(3) = 10000$
(confidential) and $\phi(4) = 15000$ (sovereign); any other $\lambda$ is
rejected with \code{ErrInvalidPrivacyLevel}. We reuse the function and its
domain unchanged.
\end{definition}

\begin{definition}[Ledger accounts]
Four maps, all held in chain state:
$\mathrm{esc}: \Buyers \to \mathbb{N}$, escrowed credit;
$\mathrm{lock}: \iota \to \mathbb{N}$, credit locked against an open order;
$\mathrm{bond}: \Ops \to \mathbb{N}$, bonded stake;
$\mathrm{earn}: \Ops \to \mathbb{N}$, claimable earnings.
\end{definition}

\begin{definition}[Epoch]
$e$, a fixed count of blocks. $A(e)$ is the multiset of claims admitted in
epoch $e$; $A_o(e)$ those attributed to operator $o$. $\Lambda(e)$ is the
emission cap for the epoch, a governance parameter.
\end{definition}

%=====================================================================
\section{The ledger}
\label{sec:ledger}

\subsection{Two flows, kept apart}

Money reaches an operator by two routes with different sources, different
bounds, and different failure modes.

\begin{description}[leftmargin=1.6em,style=nextline]
\item[Revenue.] A buyer's escrowed credit pays for work that buyer ordered. The
source is the buyer's balance. It is a transfer: total supply is unchanged.
Its bound is the order --- an operator cannot be paid more than the buyer
committed.
\item[Subsidy.] The protocol emits new AI to reward attested work in aggregate.
The source is issuance. It is a mint: supply increases. Its bound must be
external to any individual claim, because a claim's own fields are chosen by
the party being paid.
\end{description}

The current implementation braids them: it computes a mint from fields of the
claim (\S\ref{sec:admission}). A mint whose magnitude is read from a field of
the claim is unbounded, because the claimant writes the field. Separating the
two flows is the first change this document makes.

\begin{invariant}[Conservation]
\label{inv:conservation}
For every epoch $e$,
\[
\sum_{o \in \Ops} \Delta\,\mathrm{earn}(o)
\;=\;
\underbrace{\sum_{b \in \Buyers} \Delta\,\mathrm{esc}(b)^{-}}_{\text{revenue, a transfer}}
\;+\;
\underbrace{M(e)}_{\text{subsidy, a mint}},
\qquad
0 \le M(e) \le \Lambda(e),
\]
where $\Delta\,\mathrm{esc}(b)^{-}$ is the total debited from buyer $b$. No
other rule in this document increases $\mathrm{earn}$.
\end{invariant}

\begin{invariant}[Order bound]
\label{inv:orderbound}
The revenue paid against order $\Ord$ never exceeds $\pi_\Ord$, and the total
revenue paid against $\Ord$ over its lifetime never exceeds $\mathrm{lock}(\iota)$.
\end{invariant}

\begin{invariant}[Bond solvency]
\label{inv:bond}
$\mathrm{bond}(o) \ge \beta_{\min}$ is a precondition of admitting any claim
attributed to $o$. A reduction of $\mathrm{bond}(o)$ never makes it negative;
a reduction is $\min(\mathrm{bond}(o), s)$ for the scheduled amount $s$.
\end{invariant}

\subsection{Credit in}

\textbf{Proposed.} Credit enters by an escrow deposit and leaves by withdrawal
or by payment. Three transitions:

\begin{align*}
\textsc{Deposit}(b, v)&: && \mathrm{esc}(b) \mathrel{+}= v, \text{ against a transfer of } v \text{ into the settlement contract} \\
\textsc{Open}(\Ord)&: && \mathrm{esc}(b) \mathrel{-}= \pi_\Ord,\quad \mathrm{lock}(\iota) \mathrel{+}= \pi_\Ord \\
\textsc{Withdraw}(b, v)&: && \mathrm{esc}(b) \mathrel{-}= v \text{ if } v \le \mathrm{esc}(b)
\end{align*}

Locked credit is not withdrawable. An order that expires unfilled at height
$\delta$ releases its lock back to the buyer:
$\mathrm{lock}(\iota) \to 0$, $\mathrm{esc}(b) \mathrel{+}= \mathrm{lock}(\iota)$.

\textbf{Remark on denomination.} Credit is denominated in AI. Buyers who wish
to hold a stable unit must do so outside the chain or through a separate
instrument; introducing an oracle into the settlement path would add a trusted
component to the one place where the chain's guarantee is supposed to be
mechanical. This is a deliberate restriction, not an oversight, and it is a
cost: it exposes buyers to AI price movement between deposit and use.

\subsection{Rewards out}

An operator's claimable balance is spent by \textsc{Claim}, which transfers
$\mathrm{earn}(o)$ to $o$'s account. \textbf{Proposed:} a withdrawal that
leaves the chain (a bridge, another chain) MUST NOT execute below the Horizon
rung (\S\ref{sec:finality}).

%=====================================================================
\section{Admitting a served-compute claim}
\label{sec:admission}

\subsection{The predicate}

\begin{definition}[Admission]
$\Admit(\Claim, \Ord, h)$, evaluated at block height $h$ with block timestamp
$t$, holds exactly when conditions A1--A7 all hold. Admission is the only
event that moves money to an operator.
\end{definition}

\begin{description}[leftmargin=2.6em,style=sameline,labelwidth=2.2em]
\item[A1] \emph{Form.} $\Claim$ parses; $|\text{workProof}| \ge 78$; the
envelope's \code{sigLen} is nonzero and within bounds.
\item[A2] \emph{Authorship.} $\textsc{VerifyMLDSA}(P, \text{claim bytes}, \sigma)$.
The signature binds this exact claim to the key $P$.
\item[A3] \emph{Device binding.} The receipt's certificate chain verifies to an
embedded vendor root at the report's own timestamp; \code{teeSig} is the leaf's
signature over the 48-byte report; and the attested device identity equals
$\mathrm{BLAKE3}(P)$.
\item[A4] \emph{Freshness and quote binding.} The report timestamp $t_r$
satisfies $t - \Delta_{\text{fresh}} \le t_r \le t + \Delta_{\text{skew}}$, and
the report nonce equals $\mathrm{BLAKE3}(\iota \,\|\, \text{chainID})$.
\item[A5] \emph{Order binding.} $\Ord$ exists, is unfilled, $h \le \delta_\Ord$,
$o = o_\Ord$ or $o_\Ord = \bot$, and the claim's model commitment equals
$\mu_\Ord$.
\item[A6] \emph{Uniqueness.} $\mathrm{workID} = \mathrm{BLAKE3}(\iota \,\|\,
\mathrm{deviceID} \,\|\, \text{chainID})$ is unspent.
\item[A7] \emph{Bond.} $\mathrm{bond}(o) \ge \beta_{\min}$ and $o$ is not jailed
at $h$.
\end{description}

\subsection{Which of these exist}

\textbf{Observed.} A2 and A3 are implemented, together, in
\code{precompile/ai/cross\_chain\_mining.go}. \code{verifyAndMintWork} performs,
in order: an ML-DSA verification over the work proof; \code{verifyDeviceBinding},
which runs \code{verifyAttestation} against the embedded root pool and then
checks $\text{attestedDevice} = \mathrm{BLAKE3}(P)$; and a third check that the
work proof's own deviceID field equals $\mathrm{BLAKE3}(P)$. The verifier is
deterministic by construction --- the chain is checked at the report's embedded
timestamp rather than the wall clock, and the root pool is explicit rather than
the host store (\code{precompile/ai/tee\_attest.go}, \code{tee\_roots.go}) ---
and it fails closed: an empty or non-terminating chain cannot verify.

A6 exists in a weaker form: \code{ComputeWorkId(deviceId, nonce, chainId)} keyed
on the work proof's nonce field, with \code{IsSpent}/\code{MarkSpent} over a
\code{BLAKE3("spnt" || workId)} slot.

A1 is implemented. A4, A5 and A7 do not exist.

\subsection{Why the gaps matter}

\begin{observation}[The reward magnitude is written by the party being paid]
\label{obs:magnitude}
\textbf{Observed.} \code{CalculateReward} reads \code{privacy} and
\code{computeMinutes} directly out of the work proof and returns
\[
\texttt{reward} = 10^{18} \cdot \texttt{computeMinutes} \cdot
\frac{\texttt{multiplier}}{10000},
\]
with $\texttt{multiplier} \in \{2500, 5000, 10000, 15000\}$
(\code{privacyMultipliers}). \code{mintWork} calls it after the spent check and
returns the amount for the caller to mint. \textbf{Derived.} Neither field is
covered by the TEE report --- \code{verifyDeviceBinding} constrains only the
attested device identity and the signing key --- so both are chosen freely by
whoever holds the ML-DSA key. \code{computeMinutes} is a \code{uint32}, so its
maximum is $4{,}294{,}967{,}295$; at the \code{PrivacySovereign} multiplier the
single-claim maximum is $6{,}442{,}450{,}942.5$~AI. That exceeds the total
supply stated in HIP-0001 by more than six times.
\end{observation}

\begin{observation}[The replay key is also written by the claimant]
\label{obs:replay}
\textbf{Observed.} The spent-set key is
$\mathrm{BLAKE3}(\text{deviceID} \| \text{nonce} \| \text{chainID})$, and
\code{nonce} is \code{workProof[32:64]}, a field of the claim. The report's own
nonce, \code{receipt[40:48]}, is never compared to it. \textbf{Derived.} One
attested device can produce an unbounded family of distinct, individually valid
claims by varying that field; the spent set bounds replay of an identical claim,
not issuance of fresh ones. Combined with
Observation~\ref{obs:magnitude}, a single genuine TEE device is sufficient to
mint without bound.
\end{observation}

\begin{observation}[One quote serves forever]
\label{obs:staleness}
\textbf{Observed.} \code{verifyAttestation} checks the certificate chain's
validity \emph{at the report's own timestamp}, which is what makes it
deterministic, and nothing compares that timestamp to the block. \textbf{Derived.}
A quote produced once remains admissible indefinitely, and remains admissible
after the device has been decommissioned or its key exfiltrated. A4 is the fix:
the block-relative window is the liveness check the deterministic verifier
deliberately does not perform, and it belongs at the admission site, where a
deterministic block timestamp is available --- the attestation precompile
already threads exactly that value
(\code{accessibleState.GetBlockContext().Timestamp()}) and states the rule
``never \code{time.Now()}''.
\end{observation}

\begin{observation}[The compute market's verification is a tautology]
\label{obs:tautology}
\textbf{Observed.} \code{precompile/compute} at
\code{0x0300\ldots0010} exposes \code{VerifyCompute(jobID, attestation)}, whose
body recomputes $\mathrm{keccak256}(\text{jobID} \| \text{outputHash})$ from its
own storage and compares it to the supplied \code{attestation}
(\code{compute/contract.go}). \textbf{Derived.} The comparand is a function of
values the caller already supplied through \code{ClaimReward}; any caller can
compute it. The check establishes nothing about a TEE, a device, or an
execution. \code{ClaimReward} itself requires only that the caller's provider
nonce is nonzero --- that is, that the caller once called
\code{RegisterProvider} --- and transfers no value. \textbf{Specification.} A
claim MUST be admitted by $\Admit$; \code{compute}'s
\code{VerifyCompute} MUST NOT be used as an admission decision. The two real
verifiers in the block --- \code{ai.VerifyTEE} and the
\code{attestation} adapter over it, which the latter's package documentation
already names as ``exactly one way to verify a TEE quote across the block'' ---
are the only ones a settlement decision may consult.
\end{observation}

\subsection{Settlement}

\begin{algorithm}[H]
\caption{Settle a claim. Executes only if $\Admit(\Claim,\Ord,h)$ holds.}
\begin{algorithmic}[1]
\State $u \gets \min(m_\Claim,\ \bar m_\Ord)$ \Comment{meter bounded by the order}
\State $\rho \gets \min(\pi_\Ord,\ \mathrm{tariff}(\mu_\Ord)\cdot u)$ \Comment{revenue}
\State $\mathrm{lock}(\iota) \mathrel{-}= \rho$; \quad $\mathrm{earn}(o) \mathrel{+}= \rho$
\State $\mathrm{esc}(b) \mathrel{+}= \mathrm{lock}(\iota)$; \quad $\mathrm{lock}(\iota) \gets 0$ \Comment{refund the unspent cap}
\State $W_o(e) \mathrel{+}= \phi(\lambda_\Claim)\cdot u$ \Comment{accrue the epoch weight, \S\ref{sec:reward}}
\State $\textsc{MarkSpent}(\mathrm{workID})$
\State mark $\Ord$ filled; record $(y_\Claim,\ o,\ h)$
\end{algorithmic}
\end{algorithm}

Steps 1--4 are the whole of the revenue flow, and they touch no issuance. Step 5
is the operator's only influence on the subsidy flow, and it is an addend to a
share of a fixed pool, not an amount.

\begin{proposition}[Revenue is bounded by the buyer's own commitment]
Under A5 and A6, the total revenue an operator can obtain from order $\Ord$ is
at most $\pi_\Ord$.
\end{proposition}
\begin{proof}
A5 requires the order to be unfilled and step 7 marks it filled, so at most one
claim settles per order. Step 2 caps $\rho$ at $\pi_\Ord$, and step 3 debits
$\mathrm{lock}(\iota) = \pi_\Ord$, which \textsc{Open} funded from
$\mathrm{esc}(b)$. A6 additionally prevents two claims from sharing a workID
within one order, since $\iota$ is in the key.
\end{proof}

\begin{remark}[Why A6 changes the key]
Deriving the workID from $\iota$ rather than a claim-chosen nonce converts the
spent set from a duplicate filter into a one-claim-per-order rule. The change is
one line at the call site of \code{ComputeWorkId}; the spent-set machinery,
including its \code{BLAKE3} slot derivation and its atomic
check-compute-mark ordering (which the current code correctly notes avoids a
check-then-act race), is reused unchanged.
\end{remark}

\subsection{The meter}
\label{sec:meter}

\textbf{Proposed.} The meter unit is the same quantity the off-chain plane
already records per request, so that the number the chain pays on and the number
the invoice shows have one definition. \textbf{Observed:} the off-chain
recorder's wire type is deliberately narrow --- subject, namespace, an exact
decimal amount, currency, model and provider --- and carries no prompt content
and no PII (HIP-1313). The chain reads $m$, the meter reading, from the claim
and clamps it to $\bar m_\Ord$; the buyer set $\bar m_\Ord$ when opening the
order.

\textbf{Assumed.} The meter is not attested. A TEE quote of the form verified
here binds a device identity and a key; it does not bind a token count or an
elapsed time. Clamping to the buyer's cap is therefore the whole of the
protection on the revenue side, and it is sufficient there precisely because
the buyer chose the cap. It is \emph{not} sufficient on the subsidy side, which
is why the subsidy is a pool share (\S\ref{sec:reward}) rather than a function
of $m$ alone. Binding a meter to an execution requires either deterministic
re-execution (\code{precompile/inference} makes this possible for the models it
covers) or a proof system; both are out of scope here and named as open work in
\S\ref{sec:open}.

%=====================================================================
\section{Operator bonds}
\label{sec:bond}

\subsection{Where the bond lives, and where it does not}

\begin{observation}[The P-chain holds no slashable bond for a sovereign L1]
\textbf{Observed.} \code{RegisterL1ValidatorTx} carries a \code{Balance} field
documented as the validator's funding, and \code{IncreaseL1ValidatorBalanceTx}
is documented as topping up ``an L1 validator's \emph{continuous-fee}
balance''. On removal, the executor refunds the remainder to the registered
\code{RemainingBalanceOwner} as a UTXO
(\code{txs/executor/standard\_tx\_executor.go}). \textbf{Derived.} That balance
pays for the validator's continued presence; it is refundable and is not a bond
against misbehaviour. \textbf{Observed.} The consensus package that detects
equivocation states the division explicitly: ``Actual stake deduction happens in
the platformvm; this package provides detection and evidence only''
(\code{consensus/core/slashing/slashing.go}). \textbf{Observed.} We found no
stake-deduction path in \code{vms/platformvm/txs/executor}.
\end{observation}

\textbf{Specification.} The bond is held by the Hanzo staking contract --- the
manager named in \code{ConvertNetworkTx} --- on the Hanzo EVM. The Lux P-chain
holds a \emph{mirror} of the validator set's weights, not the value behind them.
This is the direct consequence of choosing \code{Manager: Contract}, and it is
the correct place for the bond: the faults that matter here are settlement
faults, and settlement state is on the Hanzo EVM.

\subsection{Weight is a projection of the bond}

\textbf{Observed.} The manager changes a validator's weight by sending a Warp
\code{L1ValidatorWeight} message --- $(\text{validationID}, \text{nonce},
\text{weight})$, 49 bytes --- which the P-chain applies through
\code{SetL1ValidatorWeightTx}. The executor verifies the message's nonce is at
least the validator's \code{MinNonce}, and verifies via \code{verifyL1Conversion}
that the message came from the L1's registered manager chain and manager
address. Two behaviours are load-bearing:

\begin{enumerate}[leftmargin=2.2em]
\item \code{weight == 0} is \emph{removal}, not suspension: the executor
refunds the remaining fee balance and deletes the validator, and refuses if this
would remove the last validator (\code{errRemovingLastValidator}).
\item \code{nonce == MaxUint64} is reserved for removal and is rejected with any
nonzero weight (\code{L1ValidatorWeight.Verify}, \code{ErrNonceReservedForRemoval}).
\end{enumerate}

\textbf{Specification.} Weight is a monotone function of the bond,
$w_o = \lfloor \mathrm{bond}(o) / \tau_{\min} \rfloor$, recomputed by the manager
whenever the bond changes and pushed as a weight message. \textbf{Proposed:}
jailing MUST be expressed as exclusion inside the manager and, on the P-chain,
as a weight \emph{reduction} to a positive floor --- never as weight zero, which
is irreversible and requires re-registration with a fresh
\code{RegisterL1ValidatorTx} and proof of possession. A design that jails by
sending weight zero silently converts a temporary penalty into permanent
ejection.

\subsection{Reading stake from inside the EVM}

\textbf{Observed.} \code{evm/precompile/contracts/stakeweight} at
\code{0x0200\ldots0006} lets a contract read a validator's P-chain weight and
the set's total weight at a height, via a BLS-signed \code{Ballot} (168 bytes:
nodeID, voter, authEpoch, P-chain height, weight, total weight, signature). The
signed statement separates chain, node, voter and epoch, so one authorisation
cannot be replayed onto another chain, node, address or revoked epoch; height
and weight are deliberately \emph{not} signed and are checked against P-chain
state instead, which is what lets one authorisation serve every later snapshot.
The expensive work is paid in \code{PredicateGas} before execution; the
in-execution read is 2{,}000 gas.

\textbf{Specification.} The settlement contract reads participation weight
through this precompile rather than trusting a self-reported figure. It
therefore depends on \code{stakeWeightConfig} being activated, which it is not
today (\S\ref{sec:chain}).

\subsection{Faults and the schedule}

\begin{description}[leftmargin=2.6em,style=sameline,labelwidth=2.2em]
\item[F1] \emph{Non-delivery.} An order assigned to $o$ reaches its deadline
$\delta$ with no admitted claim. Anyone may submit the evidence, which is the
order identifier and the height; the chain verifies the absence itself from its
own state. Reduction $s_1 = \min(\mathrm{bond}(o), \kappa_1 \pi_\Ord)$.
\item[F2] \emph{Invalid attestation.} A claim attributed to $o$ satisfies A2 ---
so authorship is attributable to a key bound to $o$ --- and fails A3 or A4. The
evidence is the claim itself. Reduction
$s_2 = \min(\mathrm{bond}(o), \kappa_2 \pi_\Ord)$, rate-limited to at most
$n_2$ occurrences per epoch per operator.
\item[F3] \emph{Output disagreement.} For an order whose model commitment names
a deterministic profile, a challenger re-executes and submits an output
commitment $y' \ne y$ together with the re-execution witness; the chain
adjudicates by the deterministic path. Reduction
$s_3 = \min(\mathrm{bond}(o), \kappa_3 \pi_\Ord)$ with $\kappa_3 > 1$.
\item[F4] \emph{Consensus fault.} Double-vote, double-sign or downtime, as
classified by \code{consensus/core/slashing} (\code{DoubleVote}, \code{DoubleSign},
\code{Downtime}, with an \code{Evidence} record carrying the serialized
conflicting votes and a \code{SlashRecord} carrying a jail expiry). Reduction
$s_4$, a fixed fraction of the bond.
\end{description}

\begin{parameter}[Slash coefficients]
$\kappa_1, \kappa_2, \kappa_3, n_2, s_4, \beta_{\min}, \tau_{\min},
\Delta_{\text{fresh}}, \Delta_{\text{skew}}$ are governance parameters. This
document fixes their \emph{constraints}, not their values.
\end{parameter}

\begin{invariant}[Deterrence]
\label{inv:deterrence}
$\kappa_1 \ge 1$: the penalty for accepting an order and not serving it must be
at least the revenue foregone by serving it, or non-delivery is free whenever
the operator finds a better use of the device. $\kappa_3 > 1$: the penalty for a
wrong answer must exceed the payment for a right one, or fabricating output is
profitable at any nonzero detection probability below one.
\end{invariant}

\begin{remark}[F2 is about attribution, not loss]
An invalid attestation never settles --- $\Admit$ rejects it --- so the ledger
loses nothing. F2 exists because submitting invalid attestations is a cheap way
to consume verification gas and to probe the verifier, and because A2 makes the
submission attributable to a bonded party. Its reduction should be small and
rate-limited; treating it like F3 would punish a misconfigured operator as
harshly as a lying one.
\end{remark}

\begin{remark}[What F1 can and cannot see]
F1 is decidable from chain state alone: the order exists, the height has passed,
no claim was admitted. It does not distinguish an operator who refused from one
whose network failed. That is a deliberate limitation of a mechanism that must
be adjudicated by a chain with no view of the operator's environment, and it
argues for $\kappa_1$ near its lower bound rather than far above it.
\end{remark}

%=====================================================================
\section{Distributing reward for served, attested compute}
\label{sec:reward}

\subsection{Pool share, not per-claim mint}

\textbf{Proposed.} At the close of epoch $e$, each operator's subsidy is its
share of a fixed pool:

\[
S_o(e) \;=\; \left\lfloor \Lambda(e) \cdot \frac{W_o(e)}{W(e)} \right\rfloor,
\qquad
W_o(e) = \sum_{\Claim \in A_o(e)} \phi(\lambda_\Claim)\, u_\Claim,
\qquad
W(e) = \sum_{o \in \Ops} W_o(e),
\]

with $u_\Claim$ the order-clamped meter of \S\ref{sec:meter} and $\phi$ the
privacy multiplier of \S\ref{sec:objects}. If $W(e) = 0$ then $S_o(e) = 0$ for
all $o$ and $M(e) = 0$: an epoch with no admitted work mints nothing. The
remainder $\Lambda(e) - \sum_o S_o(e)$, at most $|\Ops| - 1$ units lost to the
floors, is not minted.

The multiplier values and the basis-point arithmetic are the deployed ones. What
changes is the quantity they multiply --- the order-clamped meter rather than
the claim's own \code{computeMinutes} --- and the fact that the product is a
weight rather than an amount.

\begin{proposition}[The mint is bounded regardless of claim contents]
$M(e) = \sum_o S_o(e) \le \Lambda(e)$ for every epoch, for every set of
admitted claims, whatever values the claims carry in their meter and privacy
fields.
\end{proposition}
\begin{proof}
Each $S_o(e)$ is a floor of $\Lambda(e) \cdot W_o(e)/W(e)$ with
$W_o(e) \ge 0$ and $\sum_o W_o(e) = W(e)$; the unfloored terms sum to exactly
$\Lambda(e)$ when $W(e) > 0$, and flooring only decreases them. The claim's
fields influence the \emph{ratio} $W_o/W$ and therefore the split between
operators, never the total.
\end{proof}

\begin{remark}[What a pool share does not fix]
It bounds inflation; it does not make the split honest. An operator that
inflates $u$ takes share from operators that do not, and the clamp
$u \le \bar m_\Ord$ bounds that inflation only to what some buyer was willing to
pay for. Under A5 an inflated $u$ therefore costs the inflater real escrowed
credit, which is the mechanism that makes the share approximately truthful ---
not the attestation. We state this plainly because it is the load-bearing
economic assumption of the design and it is not proved here: it requires that
a self-dealing buyer-operator pair cannot recycle credit at a profit, which
holds only if the revenue leg is a genuine transfer with a nonzero cost
(fees, or the opportunity cost of locked credit). See \S\ref{sec:open}.
\end{remark}

\subsection{Where the emission schedule comes from}
\label{sec:disagreement}

\textbf{Observed.} The published record disagrees. HIP-0001 states a total
supply of $1{,}000{,}000{,}000$~AI with an emission of $100$M in year one
halving every four years to a terminal $6.25$M. The \code{hanzo-ai-chain} paper
in this repository states a hard maximum nominal supply of two trillion, split
one trillion to the DAO and one trillion to proof-earned rewards on a
Bitcoin-shaped geometric halving. HIP-0024 specifies a base staking reward of
8\% APY with a 2--5\% AI-mining bonus, and names the consensus engine
\code{nova}, which is a rung of the Lux finality ladder rather than an engine
selector.

\textbf{Specification.} This document does not resolve that disagreement. It
requires only that $\Lambda: e \to \mathbb{N}$ is a total function fixed in
chain state before epoch $e$ opens, that it is monotone non-increasing in $e$
unless changed by governance, and that $\sum_e \Lambda(e)$ converges if a supply
cap is claimed. The settlement rules of \S\ref{sec:ledger} and
\S\ref{sec:reward} hold for any $\Lambda$ satisfying these conditions, so
resolving the disagreement is a governance question and not a precondition for
implementing them.

%=====================================================================
\section{The on-chain / off-chain split}
\label{sec:split}

\subsection{The rule}

\begin{quote}
The chain holds money and commitments. It never holds content. A number the
chain pays on MUST be derivable from an on-chain order and a verified
attestation; it MUST NOT come from an off-chain assertion alone.
\end{quote}

\subsection{The split}

\begin{center}
\begin{tabularx}{\textwidth}{@{}lXl@{}}
\toprule
State & Where it lives & Why \\
\midrule
Escrow, lock, bond, earnings & Hanzo EVM, settlement contract & it is money \\
Order (buyer, $\mu$, $\pi$, $\bar m$, $\delta$) & Hanzo EVM & the bound on payment \\
Admission decision & Hanzo EVM, calling \code{0x0300\ldots} & it moves money \\
Spent set & precompile state & replay is a consensus property \\
Weight mirror & Lux P-chain, via Warp & the validator set is Lux's registry \\
Model commitment & \code{modelregistry} & a 32-byte identifier \\
Slash evidence and record & Hanzo EVM & adjudication must be public \\
\midrule
Request and response payloads & off chain, never committed & content, PII \\
Model weights & off chain; only $\mu$ on chain & size \\
Per-request meter detail & \code{hanzo.cloud\_usage} warehouse & volume \\
Invoices, balances, tax & commerce ledger (HIP-1313) & not consensus \\
Pay-per-request rail & x402 v2 (HIP-1163) & an HTTP rail, not a chain rail \\
Matching and scheduling & clearing plane & an optimisation, not a ledger \\
TEE quote generation & the operator's device & by construction \\
Fleet inventory, kubeconfigs & KMS-sealed, per org (HIP-0121) & secrets \\
\bottomrule
\end{tabularx}
\end{center}

\subsection{The one bridge between the halves}

\textbf{Observed.} \code{precompile/anchor} at \code{0x0700\ldots0010} stores
$(\text{appID}, \text{height}, \text{root})$ with per-appID height monotonicity
and a 25{,}400-gas \code{submit}; \code{get}, \code{getLatest} and a
range-capped \code{list} read it back.

\textbf{Specification.} The off-chain usage ledger commits to the chain by
anchoring the Merkle root of each epoch's records under a fixed appID. This gives
a buyer a way to prove that a detail record was in the set the operator
committed to, and gives a dispute a fixed reference. \textbf{It does not make the
off-chain records authoritative for payment}: payment is decided by
$\Admit$ against the order, and the anchor is evidence for reconciliation, not
an input to settlement. An anchor whose root were an input to settlement would
be a trusted oracle, since the party that computes the root is the party being
paid; keeping it out of the payment path is what makes it a commitment instead.

%=====================================================================
\section{Binding money movements to finality}
\label{sec:finality}

\begin{center}
\begin{tabular}{@{}lll@{}}
\toprule
Movement & Required rung & Reason \\
\midrule
\textsc{Deposit}, \textsc{Open}, lock release & Nova & local, reversible with the block \\
Admission and $\mathrm{earn}$ credit & Quasar & a counterparty will act on it \\
Bond reduction (F1--F4) & Quasar & it is adverse to a third party \\
Weight message to the P-chain & Quasar & it is an export \\
Operator withdrawal off chain & Horizon & irreversible \\
Cross-chain claim of earnings & Horizon & irreversible \\
\bottomrule
\end{tabular}
\end{center}

\textbf{Derived.} The middle rows follow from the ladder's own statement that a
block below Quasar can never be exported as deterministically final to a bridge,
another chain, or an irreversible settlement, since Nova is reorgable by a lone
equivocator holding a bare majority. The last two rows follow from
\code{AuthorizesIrreversibleSettlement}, which requires the post-quantum seal.

\begin{remark}
Admission is placed at Quasar rather than Horizon deliberately. Admission
credits an internal balance; it does not release value from the chain. Requiring
the post-quantum seal for every claim would push settlement latency to the
sealing cadence for no security gain, because the balance cannot leave without
crossing the Horizon requirement anyway.
\end{remark}

%=====================================================================
\section{Threat model}
\label{sec:threat}

\subsection{Assets}

Escrowed buyer credit; operator bonds; the issuance right (the ability to
increase supply); the integrity of the validator set; and buyer request content,
which the chain never holds.

\subsection{Adversaries}

\begin{description}[leftmargin=1.6em,style=nextline]
\item[A device holder.] Controls a genuine, vendor-certified TEE and its
attestation key. Can produce arbitrarily many valid quotes for that device.
\emph{This is the adversary the current implementation does not survive}
(Observations~\ref{obs:magnitude}, \ref{obs:replay}).
\item[A registered operator.] Holds a bond and serves orders. May under-serve,
fabricate output, or inflate the meter.
\item[A self-dealing pair.] A buyer and an operator under one principal,
recycling credit to farm subsidy share.
\item[A validator minority.] Below the Quasar floor. May equivocate; cannot
export.
\item[A network adversary.] Delays or reorders messages; cannot forge a
certificate chain or an ML-DSA signature.
\end{description}

\subsection{Assumed}

The embedded vendor roots are the correct trust anchors and their private keys
are not compromised. ML-DSA at the deployed parameter set is unforgeable. A TEE
that verifies has not had its attestation key extracted. Fewer than one third of
stake is Byzantine. None of these is proved here; the first is the weakest,
because a single vendor root compromise admits arbitrary claims, and it is why
A7 (a bond) is a precondition of admission rather than an alternative to
attestation.

\subsection{Explicitly not covered}

Confidentiality of the request payload against the operator; the operator sees
it by construction unless the model runs inside the TEE, which the attestation
verified here does not establish. Correctness of the model weights behind $\mu$.
The economics of the clearing that assigns orders.

%=====================================================================
\section{Activation}
\label{sec:activation}

\textbf{Proposed.} The changes this document requires, in dependency order:

\begin{enumerate}[leftmargin=2.2em]
\item Activate \code{stakeWeightConfig} in the Hanzo chain's upgrade
configuration; the module is registered
(\code{evm/precompile/registry/registry.go}) and implemented but absent from the
Lux mainnet \code{upgrade.json}.
\item Add the four missing admission conditions. A4 needs the block timestamp
already threaded into the attestation precompile and a nonce comparison; A5,
A6 and A7 need the order and bond state that the settlement contract holds. A6
is a change of argument to \code{ComputeWorkId}, not new machinery.
\item Deploy the settlement contract and register it as the manager address in
\code{ConvertNetworkTx}; until then the L1 cannot be sovereign with
\code{Manager: Contract}.
\item Change the subsidy path from a per-claim mint to an epoch pool share.
\code{CalculateReward} remains useful as the weight function $\phi \cdot u$;
its output stops being an amount and becomes an addend.
\item Retire \code{compute.VerifyCompute} as an admission decision
(Observation~\ref{obs:tautology}). The market records in \code{precompile/compute}
may remain as an index; they must not gate money.
\end{enumerate}

Each of steps 2, 4 and 5 is a behaviour change to an activated mainnet
precompile and therefore requires a new activation timestamp rather than an
in-place edit.

%=====================================================================
\section{Limitations and open questions}
\label{sec:open}

\begin{enumerate}[leftmargin=2.2em]
\item \emph{The meter is unattested.} Nothing in the verified quote binds a
token count or an elapsed time. The revenue side is protected by the buyer's cap;
the subsidy side is protected only by the cost of the revenue leg. Closing this
properly requires deterministic re-execution --- \code{precompile/inference} is
a byte-exact integer transformer built for exactly this, with the
CPU/Metal/CUDA/HIP parity treated as a release conformance test --- or a
succinct proof. Neither is specified here.
\item \emph{Self-dealing is bounded but not eliminated.} We assert that a
buyer-operator pair pays a real cost to recycle credit; we do not derive the
condition under which subsidy share exceeds that cost. This needs the fee and
$\Lambda$ parameters to be fixed and then measured, not argued.
\item \emph{Only one vendor root is embedded.} \code{tee\_roots.go} contains the
Intel SGX Provisioning Certification Root CA and states that chains presenting
another vendor's leaf fail closed until that vendor's root is appended.
HIP-0024's validator requirements name NVTrust GPU attestation. \textbf{Derived:}
GPU quotes do not verify today. Adding a root is a consensus change and must go
through activation.
\item \emph{F3 needs an adjudication procedure.} We give the fault and the
penalty constraint; the challenge protocol --- who may challenge, over what
window, at what deposit, and how the re-execution is bounded in gas --- is not
specified.
\item \emph{Jail semantics on the P-chain.} We specify jail as a weight
reduction to a positive floor because weight zero is removal. The interaction
between a jailed operator's reduced weight and \code{errRemovingLastValidator}
in a small validator set is not analysed.
\item \emph{Credit is denominated in AI.} Buyers carry price exposure between
deposit and use. A stable-denominated credit requires a price input in the
settlement path, which we declined; whether that trade is right is an open
question and depends on buyer behaviour we have not measured.
\end{enumerate}

%=====================================================================
\section{Summary of evidence}

File paths below are given by their leaf; the full paths are in the references.

\begin{center}
\small
\begin{tabularx}{\textwidth}{@{}Xll@{}}
\toprule
Claim & Source & Status \\
\midrule
L1 promotion is \code{ConvertNetworkTx}, own set required & \code{convert\_network\_tx.go} & observed \\
Security mode has two orthogonal axes & \code{security/mode.go} & observed \\
Weight zero is removal, not suspension & \code{standard\_tx\_executor.go} & observed \\
P-chain validator balance is a fee balance & \code{increase\_\ldots\_balance\_tx.go} & observed \\
Consensus detects faults, does not deduct stake & \code{core/slashing} & observed \\
Export needs Quasar; irreversibility needs Horizon & \code{finality.go} & observed \\
TEE verification is real and fails closed & \code{ai/tee\_attest.go} & observed \\
Device binding is $\mathrm{deviceID} = \mathrm{BLAKE3}(P)$ & \code{cross\_chain\_mining.go} & observed \\
Only the Intel SGX root is embedded & \code{ai/tee\_roots.go} & observed \\
Reward reads attacker-chosen \code{computeMinutes} & \code{ai/ai\_mining.go} & observed \\
Single-claim maximum $6.44\times10^{9}$~AI & \code{uint32} $\times$ 1.5 & derived \\
Spent-set key uses an attacker-chosen nonce & \code{ComputeWorkId} site & observed \\
\code{compute.VerifyCompute} is caller-computable & \code{compute/contract.go} & observed \\
\code{stakeWeightConfig} is not activated on mainnet & \code{upgrade.json} & observed \\
Admission predicate A1--A7 & this document & proposed \\
Two-flow ledger and pool-share subsidy & this document & proposed \\
Fault classes F1--F4 and their constraints & this document & proposed \\
\bottomrule
\end{tabularx}
\end{center}

%=====================================================================
\section*{References}
\addcontentsline{toc}{section}{References}

\begin{enumerate}[label={[\arabic*]},leftmargin=2.6em]
\item Lux P-chain transaction set. \code{luxfi/node}, \code{vms/platformvm/txs}:
\code{create\_network\_tx.go}, \code{create\_chain\_tx.go},
\code{add\_chain\_validator\_tx.go} (deprecation notice),
\code{convert\_network\_tx.go}, \code{register\_l1\_validator\_tx.go},
\code{set\_l1\_validator\_weight\_tx.go},
\code{increase\_l1\_validator\_balance\_tx.go}, \code{disable\_l1\_validator\_tx.go}.
\item Lux network security model. \code{luxfi/node},
\code{vms/platformvm/security/mode.go}; LP-018, network security modes.
\item Lux Warp validator messages. \code{luxfi/node},
\code{vms/platformvm/warp/message}: \code{chain\_to\_l1\_conversion.go},
\code{l1\_validator\_registration.go}, \code{l1\_validator\_weight.go}.
\item Lux finality ladder. \code{luxfi/consensus},
\code{engine/chain/finality.go}, \code{engine/chain/cert.go},
\code{config} stake floors, \code{conformance/corpus.json}.
\item Equivocation detection. \code{luxfi/consensus}, \code{core/slashing}.
\item AI mining precompile. \code{luxfi/precompile}, \code{ai}:
\code{ai\_mining.go}, \code{tee\_attest.go}, \code{tee\_roots.go},
\code{cross\_chain\_mining.go}.
\item Attestation precompile. \code{luxfi/precompile}, \code{attestation}.
\item Compute market precompile. \code{luxfi/precompile}, \code{compute}.
\item Model registry and deterministic inference. \code{luxfi/precompile},
\code{modelregistry}, \code{inference}.
\item Anchor precompile. \code{luxfi/precompile}, \code{anchor}.
\item Stake-weight precompile. \code{luxfi/evm},
\code{precompile/contracts/stakeweight}.
\item Post-quantum precompile block, slots \code{0x012201}--\code{0x012208}.
\code{luxfi/precompile}: \code{mlkem}, \code{mldsa}, \code{slhdsa}, \code{pulsar},
\code{p3q}, \code{corona}, \code{magnetar}, \code{hqc}; LP-4200.
\item Mainnet precompile activation. \code{luxfi/genesis},
\code{configs/mainnet/upgrade.json}.
\item HIP-0001, AI --- Hanzo's native currency.
\item HIP-0024, Hanzo Sovereign L1 Chain Architecture.
\item HIP-0096, AI Compute Contribution Rewards.
\item HIP-0121, BYO Compute Fleet \& Metered Billing.
\item HIP-1020, Chain Registry.
\item HIP-1163, x402 --- Pay Per Request.
\item HIP-1313, Usage --- The Metered Record.
\item \emph{Clearing Heterogeneous Compute}, Hanzo AI Research, this repository.
\item \emph{Hanzo AI Chain (AIVM)}, Hanzo AI Research, this repository.
\end{enumerate}

\end{document}
